\documentclass{article}
\usepackage{amsmath}
\usepackage{braket}
\usepackage{mathtools}
\DeclarePairedDelimiter\abs{\lvert}{\rvert}%
\newcommand{\ketbra}[2]{|#1\rangle\langle#2|}

\begin{document}
\section*{Crack or no crack?}

In order to apply the methods of quantum hypothesis testing, we must first describe the family of states corresponding to surface cracks with different widths $\delta x$ and depths $\delta z$. This can be done using an incoherent ensemble of three point sources

\begin{equation}
    \rho(\delta x, \delta z) = \frac{\ketbra{\psi_1}{\psi_1} + \ketbra{\psi_2}{\psi_2} + \ketbra{\psi_3}{\psi_3}}{3}
\end{equation}

where
\begin{equation}
    \begin{aligned}
        \psi_1 & = \psi(x\!-\!\frac{\delta x}{2},y,0)\\
        \psi_2 & = \psi(x\!+\!\frac{\delta x}{2},y,0)\\
        \psi_3 & = \psi(x,y,-\delta z).
    \end{aligned}
\end{equation}

We can then consider a binary hypothesis test between the existence of a crack and it's absence.
Importantly, the null hypothesis represents the case of no crack described by our model as three collinear sources with corresponding density matrix $\rho(\delta x, 0)$.
Whereas the alternative hypothesis represents a crack with a specific depth $\delta z$ given by the density matrix $\rho(\delta x, \delta p)$.

From these density matrices we then numerically calculate the quantum Chernoff bound.
However, this first requires us to represent each state in an orthonormal basis. To do this, we define an additional point source

\begin{equation}
    \psi_4 = \psi(x,y,0)
\end{equation}

and consider the Gram matrix $G_{u,v} = \braket{\psi_u}{\psi_v}$ between the states $\psi_1, \psi_2, \psi_3$ and $\psi_4$.
To find a set of vectors with overlaps given by $G$, which is called a vector realisation of $G$, we can use the fact that the matrix is positive semidefinite to decompose it into $G = A^\dagger A$.
Notably, this gives many different choices for $A$ however this does not affect the quantum Chernoff bound as it is invariant under change of basis. For simplicity, we choose $A = \sqrt{G}$ which can be calculated numerically by finding the eigensystem of $G$.

Once we have found a vector realisation, we construct the corresponding density matrices of $\rho(\delta x,0) $ and $\rho(\delta x, \delta z)$.
The eigenvalues $\lambda_i$, $\mu_j$ and eigenvectors $\ket{u_i}$, $\ket{v_j}$ can be used to calculate

\begin{equation}
    \mathrm{Tr}\left[\rho(\delta x,0)^s \rho(\delta x, \delta z)^{1-s} \right] = \sum_{ij} \lambda_i^s \mu_j^{1-s} \abs{\braket{u_i|v_j}}^2.
\end{equation}

Minimisation over $s$ can then be performed to find

\begin{equation}
\xi_Q = -\log \min_{0 \leq s \leq 1} \mathrm{Tr} \left[ \rho(\delta x,0)^s \rho(\delta x, \delta z)^{1-s} \right].
\end{equation}




As well as numerically calculating the quantum Chernoff bound, we also found an upper bound to the lowest order of $\delta z$. To do this we use the fact that the QCB can be bounded by the quantum fidelity which is defined as

\begin{equation}
F(\rho_0, \rho_1) = \mathrm{Tr}\left( \sqrt{\sqrt{\rho_0} \rho_1 \sqrt{\rho_0}}\right).
\end{equation}

Using results from [QCB paper] we obtain bounds on the QCB

\begin{equation}
- \log F(\rho_0, \rho_1)  \leq \xi_Q \leq - 2 \log F(\rho_0, \rho_1).
\end{equation}

To calculate a low order expansion of the fidelity we use a method from [Lupo Paper] that calculates the fidelity from purifications of $\rho_0$ and $\rho_1$.
The state $\rho$ can be represented by an ensemble of pure states

\begin{equation}
\rho = \sum_j p_j \ketbra{\psi_j}{\psi_j}
\end{equation}

with pure state $\ket{\psi_j}$ defined using an orthonormal basis $\{\ket{i}\}$

\begin{equation}
\ket{\psi_j} = \sum_i \braket{i|\psi_j} \ket{i}.
\end{equation}

We can then write the purification of $\rho$ as

\begin{equation}
\ket{\Psi} = \sum_{ij} C_{ij} \ket{i} \ket{j}
\end{equation}

where $C_{ij} = \sqrt{p_j} \braket{i|\psi_j}$ and $\{\ket{j}\}$ is an orthonormal basis representing each source.

The fidelity is then given by the formula

\begin{equation}
F(\rho_0, \rho_1) = \max_V \abs{\braket{\Psi_0| I \otimes V | \Psi_1}}
\end{equation}

where $\ket{\Psi_0}$ and $\ket{\Psi_1}$ are the purifications of $\rho_0$ and $\rho_1$. By representing each purification using the matrices $C^{(0)}$ and $C^{(1)}$ a simple expression for the fidelity can be obtained

\begin{equation}
\begin{aligned}
F(\rho_0, \rho_1) =& \max_V \left| \sum_{jk} V_{jk} \sum_{i} C^{(0)^*}_{ij} C^{(1)}_{ik} \right| \\
=& \max_V \left| \text{Tr}(V^T M) \right| = \|M\|_1
\end{aligned}
\end{equation}

where $\|M\|_1 = \mathrm{Tr}(\sqrt{M^\dagger M})$ is the trace norm and $M = C^{(0)^\dagger} C^{(1)} = \sum_i C^{(0)^*}_{ij} C^{(1)}_{ik}$ is the Gram matrix calculated from the weighted pure states in $\rho_0$ and $\rho_1$

\begin{equation}
\begin{aligned}
M_{jk} =& \sum_i (C_0^*)_{ij} (C_1)_{ik}\\
= & \sum_{i} (\sqrt{p_j}\braket{i|\psi_j})^*(\sqrt{q_k} \braket{i|\phi_k})\\
= & \sqrt{p_j q_k} \sum_i \braket{\psi_j|i} \braket{i|\phi_k} = \sqrt{p_j q_k} \braket{\psi_j|\phi_k}.
\end{aligned}
\end{equation}

It is important to note that $M$ remains the invariant under changes of the computational basis $\ket{i}$.
To show this, let $\ket{i'} = U\ket{i}$ be a change of basis where $U$ is a unitary matrix. The corresponding matrix of a purified state in this basis is

\begin{equation}
C'_{ij} = \braket{i'|\psi_j} = \braket{i|U^\dagger|\psi_j}.
\end{equation}

and therefore $C' = U^\dagger C$.  Applying the same transformation to both $C^{(0)}$ and $C^{(1)}$ gives

\begin{equation}
M' = C^{(0)^\dagger} U U^\dagger C^{(1)} =   C^{(0)^\dagger} C^{(1)} =  M.
\end{equation}

Similarly, we can show that $\|M\|_1$ is invariant under change of ensemble choice. Given that $\rho$ can be represented by two different ensembles

\begin{equation}
\rho = \sum_j p_j \ketbra{\psi_j}{\psi_j} = \sum_k q_k \ketbra{\psi_k}{\psi_k}
\end{equation}

the HJW theorem shows that these different ensembles are related by a unitary transformation

\begin{equation}
\sqrt{q_k} \ket{\phi_k} = \sum_j V_{kj} \sqrt{p_j} \ket{\psi_j}.
\end{equation}

Therefore, we can show that corresponding purification matrices of each ensemble are related by this unitary transformation

\begin{equation}
\begin{aligned}
C'_{ik} = & \sqrt{q_k} \braket{i|\phi_k} \\
= & \bra{i}\left(\sum_j V_{kj} \sqrt{p_j} \ket{\psi_j} \right) \\
= & \sum_j V_{kj} \sqrt{p_j} \braket{i|\psi_j}\\
= & \sum_j V_{kj} C_{ij}
\end{aligned}
\end{equation}

and consequently $C' = CV^{T}$.
We can represent changes in ensemble for $\rho_0$ and $\rho_1$ as $C^{(0)} V^{(0)}$ and $C^{(1)} V^{(1)}$ which gives us the new matrix

\begin{equation}
M' = V^{(0)^\dagger} C^{(0)^\dagger} C^{(1)} V^{(1)} = V^{(0)^\dagger} M V^{(1)}.
\end{equation}

We can then show that $\|M'\|_1 = \|M\|_1$ using the fact that the trace norm is unitary invariant [Provide Source].

Using these properties we have demonstrated, we can define a matrix $M$ from which we can calculate the fidelity.
To do this, we introduce the symbols $\alpha$, $\beta$ and $\gamma$ to describe the overlaps between each of our point sources

\begin{equation}
\begin{aligned}
\alpha =& \braket{\psi_1|\psi_4} = \braket{\psi_2|\psi_4} = \exp \left(-\frac{\delta x^2}{16} \right) \\
\beta =& \braket{\psi_1|\psi_3} = \braket{\psi_2|\psi_3} = \frac{2i}{\delta z + 2i} \exp \left(-\frac{i \delta x^2}{8(\delta z + 2i)} - i k z_R \delta z \right) \\
\gamma =& \braket{\psi_4|\psi_3} = \frac{2i}{\delta z + 2i} \exp \left( -i k z_R \delta z \right)
\end{aligned}
\end{equation}

and additionally,

\begin{equation}
\braket{\psi_1|\psi_2} = \exp \left(-\frac{\delta x^2}{4} \right) = \alpha^4.
\end{equation}

 Furthermore , we introduce two new states to simplify the matrix $M$,

\begin{equation}
\begin{aligned}
\ket{+} =&\frac{\ket{\psi_1} + \ket{\psi_2}}{\sqrt{2 + 2\alpha^4}} \\
\ket{-} =& \frac{\ket{\psi_1} - \ket{\psi_2}}{\sqrt{2 - 2\alpha^4}}.
\end{aligned}
\end{equation}

This then gives us an alternative description of our quantum states

\begin{equation}
\begin{aligned}
\rho(\delta x, 0) =& \frac{(1+\alpha^4)\ketbra{+}{+} + (1-\alpha^4)\ketbra{-}{-} + \ketbra{\psi_4}{\psi_4}}{3} \\
\rho(\delta x, \delta z) =& \frac{(1+\alpha^4)\ketbra{+}{+} + (1-\alpha^4)\ketbra{-}{-} + \ketbra{\psi_3}{\psi_3}}{3}.
\end{aligned}
\end{equation}

Consequently, we can then calculate the matrix $M$ from the weighted overlaps of $\ket{+}, \ket{-}, \ket{\psi_3}$ and $\ket{\psi_4}$

\begin{equation}
M = \begin{bmatrix}
\frac{1+\alpha^4}{3} & 0 & \frac{\sqrt{2} \alpha}{3} \\
0 & \frac{1 - \alpha^4}{3} & 0 \\
\frac{\sqrt{2} \beta}{3} & 0 & \frac{\gamma}{3}.
\end{bmatrix}
\end{equation}

Since $M$ is of block diagonal form

\begin{equation}
M = \underbrace{\left[ \frac{1-\alpha^4}{3} \right]}_{M_1} \oplus \frac{1}{3} \underbrace{\begin{bmatrix} 1+\alpha^4 & \sqrt{2}\beta \\ \sqrt{2}\alpha & \gamma \end{bmatrix}}_{M_2},
\end{equation}

we can use the fact that $\|M\| = \|M_1\| + \|M_2\|$ to simplify our calculations.
Using a low order approximation of $\alpha$ we can find that that  $\|M_1\| \approx \frac{\delta x^2}{12}$.
Furthermore, we can use the $2 \times 2$ matrix formula for the trace norm

\begin{equation}
\|A\| = \sqrt{\mathrm{Tr}\left(A^\dagger A \right) + 2 \abs{\mathrm{det}(A)}}
\end{equation}

along with low order approximations of the other variables $\beta$ and $\gamma$ to find that $\|M_2\| \approx 1 - \frac{\delta x^2}{12} - \frac{\delta z^2}{24}$.
As a result this gives the low order approximation of the fidelity as

\begin{equation}
\begin{aligned}
F(\rho_0,\rho_1) = \|M\| \approx \frac{\delta x^2}{12} + 1 - \frac{\delta x^2}{12} - \frac{\delta z^2}{24} = 1 - \frac{\delta z^2}{24}
\end{aligned}
\end{equation}

\subsection*{Direct Imaging}

To calculate the Chernoff information for classical intensity measurement, we use the corresponding probability distribution found by summing the intensities for each point

\begin{equation}
P(x,y|\delta x,\delta z)=\frac{\abs{\psi_1}^2 + \abs{\psi_2}^2 + \abs{\psi_3}^2}{3}.
\end{equation}

The Chernoff information can then be calculated from the equation

\begin{equation}
\xi_{DI} = - \log \min_{0 \leq s \leq 1} \iint P(x,y|\delta x,0)^s P(x,y|\delta x,\delta z)^{1-s} dx dy.
\end{equation}

The behaviour of $\xi_{DI}$ of course can be illustrated numerically for different values of $\delta x$ and $\delta z$.
However, we can also perform a low order expansion to find the leading terms of the error exponent.
To do this, we expand the integrand in $\delta x$ and $\delta z$ and integrate term by term.
This gives us the approximation

\begin{equation}
\xi_{DI} \approx - \log \min_{0 \leq s \leq 1} \left( 1 - s(1-s) \frac{\delta z^4}{18}\right) = -\log \left(1 - \frac{\delta z^4}{72} \right)
\end{equation}

as the expression is minimised for $s = \frac{1}{2}$.
Moreover, we can use the fact that $-\log(1-x) \approx x$ to further simplify the expression to

\begin{equation}
\xi_{DI} \approx \frac{\delta z^4}{72}.
\end{equation}

\subsection*{Mode Sorting}

The probability of a mode $\ket{\phi_{j,l}}$ being activated for a crack of given width $\delta x$ and depth $\delta z$ can be calculated by summing the individual probabilities for each point as described in [EQ 32]. From this we can then calculate the Chernoff information numerically by calculating the probabilities then performing minimisation over the equation

\begin{equation}
\xi_{MS} = -\log \min_{0 \leq s \leq 1} \sum_{j,l} p_{j,l}(\delta x, 0)^s p_{j,l}(\delta x, \delta z)^{1-s}.
\end{equation}

To find a low order approximation of the Chernoff information we can expand each probability $p_{j,l}(\delta x, \delta z)$. Expanding each probability gives us only four modes with second order terms

\begin{equation}
	\begin{aligned}
		p_{0,0}(\delta x, \delta z)& \approx 1 - \frac{\delta x^2}{12} - \frac{\delta z^2}{12}\\
		p_{1,0}(\delta x, \delta z)& \approx \frac{\delta x^2}{12}\\
		p_{2,0}(\delta x, \delta z)& \approx \frac{\delta z^2}{24}\\
		p_{0,2}(\delta x, \delta z)& \approx \frac{\delta z^2}{24}.
	\end{aligned}
\end{equation}

Since $p_{2,0}(\delta x, 0) = p_{2,0}(\delta x, 0) = 0$, the terms for these modes are not necessary for the calculation of  $\xi_{MS}$.
Substituting our approximation into our equation we get

\begin{equation}
\xi_{MS} \approx -\log \min_{0 \leq s \leq 1} \left(1  - \frac{\delta z^2}{12} + \frac{\delta z^2 s}{12} \right) \approx \frac{\delta z^2}{12}.
\end{equation}


\end{document}
